\section{Materials and Methods}
% Figure 1: what the study does, in the order of the paper. Every number is a macro produced by the pipeline, so the
% figure cannot drift from the results; step 11 verifies that each one exists.
\newcommand{\boxheight}{4.75cm}
\begin{figure}[!h]
	\centering
	\resizebox{\textwidth}{!}{%
		\begin{tikzpicture}[
			font=\sffamily\scriptsize,
			>=Stealth,
			node distance=0.45cm,
			basebox/.style={
				thick,
				rounded corners=4pt,
				inner sep=6pt,
				anchor=center,
				align=center
			},
			phase1/.style={basebox, draw=red!70!black, fill=red!4},
			phase2/.style={basebox, draw=blue!75!black, fill=blue!4},
			phase3/.style={basebox, draw=teal!80!black, fill=teal!4},
			phase4/.style={basebox, draw=violet!75!black, fill=violet!4},
			phase5/.style={basebox, draw=gray!70!black, fill=gray!5},
			arrowline/.style={->, ultra thick, draw=gray!65}
			]

			% STAGE 1: AUDIT OF THE PUBLISHED PIPELINE
			\node[phase1] (box1) {
				\begin{minipage}[t][\boxheight][t]{3.15cm}
					\centering
					\textbf{\footnotesize 1. Audit} \par\vspace{2pt}\hrule\vspace{7pt}
					\textbf{Object:} \\published code, re-run as released \par\vspace{4pt}
					\textbf{Leak classes:} \\re-split after augmentation; early stopping on test; in-sample importance; duplicates \par\vspace{4pt}
					\textbf{Removed one at a time:} \\$R^2$ \KarLadderRungZeroRTwoMean\ $\rightarrow$ \KarLadderRungFourRTwoMean
				\end{minipage}
			};

			% STAGE 2: LEAKAGE-FREE BENCHMARK
			\node[phase2, right=of box1] (box2) {
				\begin{minipage}[t][\boxheight][t]{3.15cm}
					\centering
					\textbf{\footnotesize 2. Benchmark} \par\vspace{2pt}\hrule\vspace{7pt}
					\textbf{Data:} \\$N = \NCompounds$ TgDHFR \par\vspace{4pt}
					\textbf{Protocol:} \\$\OuterKFolds \times \OuterRepeats$ nested CV, \TotalOuterFolds\ shared folds \par\vspace{4pt}
					\textbf{Models:} \\4 representations vs.\ 7 neural \par\vspace{4pt}
					\textbf{Inference:} \\Nadeau--Bengio, Holm, TOST
				\end{minipage}
			};

			% STAGE 3: ROBUSTNESS
			\node[phase3, right=of box2] (box3) {
				\begin{minipage}[t][\boxheight][t]{3.15cm}
					\centering
					\textbf{\footnotesize 3. Robustness} \par\vspace{2pt}\hrule\vspace{7pt}
					\textbf{Chance correlation:} \\y-randomization, 4 forms \par\vspace{4pt}
					\textbf{Chemotype shift:} \\scaffold-grouped splits (\ScaffoldNGroups\ groups) \par\vspace{4pt}
					\textbf{Curation:} \\extreme-potency records retained \par\vspace{4pt}
					\textbf{Cost:} \\wall-clock time per model
				\end{minipage}
			};

			% STAGE 4: FINDINGS
			\node[phase4, right=of box3] (box4) {
				\begin{minipage}[t][\boxheight][t]{3.15cm}
					\centering
					\textbf{\footnotesize 4. Findings} \par\vspace{2pt}\hrule\vspace{7pt}
					\textbf{Morgan ensemble:} \\$R^2 = \MorganRTwoMean$ \par\vspace{4pt}
					\textbf{Default D-MPNN:} \\$R^2 = \ChemPropDMPNNFoldRTwoMean$ \par\vspace{4pt}
					\textbf{Pretrained D-MPNN:} \\$R^2 = \ChemeleonRTwoMean$, the only equivalent model \par\vspace{4pt}
					\textbf{Reported gain:} \\accounted for by the leaks
				\end{minipage}
			};

			% STAGE 5: APPLICATION
			\node[phase5, right=of box4] (box5) {
				\begin{minipage}[t][\boxheight][t]{3.15cm}
					\centering
					\textbf{\footnotesize 5. Application} \par\vspace{2pt}\hrule\vspace{7pt}
					\textbf{Library:} \\$N = \ScreenNLibrary$ FDA-approved \par\vspace{4pt}
					\textbf{Applicability domain:} \\\ScreenNDomain\ of \ScreenNStructural\ inside \par\vspace{4pt}
					\textbf{Orthogonal check:} \\docking, redock RMSD \RedockRMSD\,\AA \par\vspace{4pt}
					\textit{Hypotheses, not leads; no selectivity assay}
				\end{minipage}
			};

			\draw[arrowline] (box1) -- (box2);
			\draw[arrowline] (box2) -- (box3);
			\draw[arrowline] (box3) -- (box4);
			\draw[arrowline] (box4) -- (box5);

		\end{tikzpicture}%
	}
	\caption{\textbf{Structure of the study.} The published pipeline is first re-executed and taken apart one leak at a
		time, which measures what each class of leak contributes to its reported $R^2$ (1). The same data are then used
		for a leakage-free benchmark in which every model is scored on identical outer folds and every comparison is
		corrected for the overlap of training sets and tested for equivalence (2), and the conclusions are checked
		against chance correlation, a harder splitting scheme, the curation decision and computational cost (3). The
		findings (4) are that a fingerprint ensemble is not outperformed by a graph network at this data volume, that
		only a pretrained network reaches equivalence with it, and that the gain reported by the original study is
		accounted for by the leaks removed in (1). The final model is applied to approved drugs as hypothesis
		generation, with each candidate placed relative to the applicability domain (5).}
	\label{fig:pipeline_flowchart}
	\vspace{-10pt}
\end{figure}


Figure~\ref{fig:pipeline_flowchart} summarizes the audit, the benchmark, the robustness analyses and the
virtual-screening application.

% ==========================================
% 2. SECTION 2.1
% ==========================================

\subsection{Bioactivity Benchmark Dataset and Curation}
\label{sec:materials_curation}

This study uses the published benchmark of \textit{Toxoplasma gondii} dihydrofolate reductase (TgDHFR) inhibitors
compiled by Karimova and Senanayake \citep{karimova2025}, taken as released, after their curation. Working on their
data rather than on a set of our own keeps the comparison with their results on one protocol and one object, so that a
difference in score cannot be attributed to a difference in curation. The primary benchmark uses no synthetic
augmentation, which leaves the measured performance attributable to the algorithms and to the empirical measurements
alone.

The curation inherited from the original study merged drug--protein interaction parameters and SMILES structures from
BindingDB with bioactivity data from ChEMBL, standardized the structures by salt stripping and structural validation,
and resolved compounds measured more than once: assays whose bioactivity varied by more than 1.5 log units were
discarded by the interquartile-range rule and the remaining values averaged into one measurement per ligand
\citep{karimova2025}. None of those steps was repeated here. The released table, file
\texttt{1-Training\_Database/V2-combined\_bd\_chmbl\_unique.csv} of the authors' repository, contains \DataNRaw\
records, and the copy used here is identical to it.

Two decisions were taken on top of it, and both are ours. The first is an exclusion rule for extreme potencies,
specified before any model was fitted: Tukey's rule on the $pIC_{50}$ distribution (fences at Q1 $-$ 1.5\,IQR $=
\OutlierFenceLow$ and Q3 $+$ 1.5\,IQR $= \OutlierFenceHigh$) places \DataNOutliers\ records outside the fences, a
classical antifolate analogue at $pIC_{50} = \OutlierHighPIC$ and a practically inactive compound at $pIC_{50} =
\OutlierLowPIC$, which leaves \DataNAfterOutliers\ entries. Because an exclusion of this kind can itself shape the
reported performance, every conclusion of this study was also re-derived in a separate execution of the complete
benchmark that keeps the two records (Section~\ref{sec:outlier_sensitivity}).

The second decision is structural deduplication, which the Organization for Economic Co-operation and Development
(OECD) principles for QSAR validation require \citep{tropsha2010, cherkasov2014, iwuala2025} and without which the
same molecule can train and test a model. All SMILES were canonicalized with RDKit, and the \DataNAfterOutliers\
entries turned out to describe only \DataNMain\ distinct structures: \DataNDuplicatesRemoved\ molecules appeared twice
under different PubChem CIDs, and the second occurrence of each was removed. The measured $pIC_{50}$ values within
each pair differed by at most \DataDuplicateMaxSpread\ units, so the target distribution is practically unchanged. The
removed records are listed in \texttt{data/V2-combined\_bd\_chmbl\_unique\_duplicates.csv}.

Two benchmarks follow from those decisions and are kept distinct throughout. The \emph{primary benchmark} comprises
the $\OutlierNMain$ unique structures that remain once the two extreme-potency records are excluded, and carries every
result reported below unless stated otherwise. The \emph{sensitivity benchmark} comprises the $\OutlierNSens$ unique
structures obtained when those two records are retained, and is used in Section~\ref{sec:outlier_sensitivity} alone.
The target is $pIC_{50}$ ($-\log_{10}(IC_{50})$) \citep{karimova2025}; the logarithmic transformation places potency
measurements on the conventional and more numerically tractable scale for regression.

% ==========================================
% 2. SECTION 2.2 (classical replication of the published interventions)
% ==========================================

\subsection{Classical Replication of the Published Interventions}
\label{sec:materials_replication}

To quantify what the interventions reported in the original study contribute, its hold-out protocol was replicated in
an isolated ablation. A 15\% hold-out partition was separated first (seed 42), and everything below was learned from the development partition only. The baseline (A) is the stacking ensemble on all \MorganNFeaturesTotal\ Morgan bits. The feature-selection variant (B) computes the permutation importance of every bit on the development partition (5 repeats, negative mean squared error) and keeps the bits whose mean importance is positive beyond floating-point noise (greater than $10^{-9}$), because bits that the model never uses have an importance that is zero up to rounding error. The augmentation variant (C) adds, to the selected features, perturbation of the training features via additive Gaussian noise: applying the two noise levels of the published code ($\sigma = 0.01$ and $\sigma = 0.001$, which that code introduces as a factor to be adjusted to the scale of the data and does not otherwise justify) generated two synthetic variants per original molecule, which triplicated the training space (to \ExpCAugSizeTrain\ rows) with the held-out molecules left untouched. Each variant is the average of five stacking ensembles that differ in the seed of their tree models.

The leak that the original selection contains was also measured on this classical replication, where it needs no deep
network. Variant B computes its permutation importance on the development partition alone. A second version of it,
identical except that the importance is computed on the complete data set as the original study does, was fitted and
scored on the same held-out molecules, and the two sets of selected bits were compared by their Jaccard index. Variant
C was treated in the same way.

% ==========================================
% 2. SECTION 2.3 (audit of the original pipeline)
% ==========================================

\subsection{Audit of the Original Pipeline}
\label{sec:materials_audit}

The comparison above treats the figures of the original study as given. To establish what produces them, its published
code was obtained and run in its original TensorFlow environment before anything was changed, which fixes the object of
the audit as the pipeline as published rather than a reimplementation of it. Because one run of a stochastic pipeline
carries its own variability, every stage was then repeated over \KarReproNSeeds\ random seeds.

The pipeline was then taken apart one mechanism at a time. Each of the four mechanisms below is a property of the
published code, verifiable by reading it; what the audit establishes is therefore not whether a mechanism is present,
but how much of the reported $R^2$ it accounts for. Starting from the published protocol (R0), four modifications were
applied in sequence, each closing one path by which information about the test molecules reaches the fitted model:
(R1) the synthetic copies produced by Gaussian augmentation are no longer re-partitioned after augmentation, so a
perturbed copy of a test molecule can no longer train the model; (R2) early stopping and model choice use an internal
split of the training partition instead of the test set; (R3) the permutation importance that drives feature selection
is computed out of fold inside the training rows instead of on the test molecules; and (R4) structures that occur more
than once are removed from both sides of the split. Every rung was run over \KarReproNSeeds\ random seeds.

Each rung alters only the mechanism it names. What follows from that alteration---how many bits the selection retains,
and which molecules remain evaluable once duplicate structures are removed---is a consequence of it and is reported as
such, not a further change to the protocol.

The ladder therefore yields two quantities that are not of the same kind, and they are reported differently. The
differences between rungs R1 to R4 are paired: those rungs share the \KarLadderNCommonTest\ molecules common to all of
them, so each difference estimates the effect of closing one path on one fixed set of test molecules, and a percentile
bootstrap over those molecules gives its confidence interval. The difference between R0 and R1 is not paired, because
the published protocol builds its own evaluation partition by re-splitting the augmented data and therefore shares no
test molecules with the rungs below it; that difference estimates the gap between the performance the published
protocol reports and the performance of the same pipeline without the re-split, and is summarized over seeds rather
than over molecules. A reference model with no feature selection of any kind was fitted on the same clean split.

% ==========================================
% 2. SECTION 2.4
% ==========================================

\subsection{Representation Spaces: Molecular Fingerprints and Conformational Descriptors}

We generated four distinct molecular representation spaces using RDKit \citep{landrum2025} to systematically map the structure-activity relationship. The structural baseline relied on Morgan fingerprints (ECFP4). We computed these with a radius of 2 and folded them into a fixed-length vector of \MorganNFeaturesTotal\ bits. Extended-connectivity fingerprints offer rapid calculation and capture a vast array of distinct local substructural features \citep{rogers2010}; chirality was not encoded. Classical 2D fingerprint-based models frequently achieve predictive parity with state-of-the-art 3D structure-based models when relying exclusively on ligand information \citep{gao2020}.

To establish a physicochemical baseline, the pipeline extracted the complete RDKit descriptor list (\RdkitTwoDNFeaturesTotal\ conventional two-dimensional topological and physicochemical descriptors). Curated 2D molecular representations demonstrate substantial robustness; recent studies confirm they maintain scaffold-independent advantages over 3D alternatives across various property prediction tasks \citep{tiwari2026}. A third hybrid space concatenated these \RdkitTwoDNFeaturesTotal\ descriptors with the \MorganNFeaturesTotal-bit Morgan fingerprints to evaluate potential synergistic effects.

A fourth representation appended ten 3D shape descriptors (asphericity, eccentricity, inertial shape factor, the two normalized principal moment ratios, the three principal moments of inertia, radius of gyration, and spherocity index) to this hybrid, for \RdkitTwoDThreeDFpNFeaturesTotal\ features in total, to assess the predictive value of spatial geometry. Hydrogens were added and a single conformer per molecule was generated with the Experimental-Torsion Knowledge Distance Geometry algorithm (ETKDGv3, fixed random seed of 42, random initial coordinates as a fallback) \citep{riniker2015}. Because purely distance-based constraints often distort aromatic rings and $sp^2$ centers, ETKDGv3 integrates experimental torsion-angle preferences derived from small-molecule crystallographic data. The conformer was then locally minimized with the MMFF94 force field before descriptor extraction. Conformer embedding failed for \RdkitTwoDThreeDFpNThreeDFailed\ of the \DataNMain\ molecules, and \RdkitTwoDThreeDFpNCellsMissing\ descriptor values of the 2D+3D+FP matrix were missing or infinite before imputation.

Although integrating spatial information remains critical for molecular image-based modeling \citep{matsu2023} and geometry-augmented representation learning \citep{yanan2024}, recent empirical evaluations suggest that 3D conformer descriptors yield limited incremental predictive value over curated 2D representations in standard tabular benchmarks \citep{tiwari2026}. Our nested cross-validation framework explicitly tested this hypothesis. The purely structural Morgan fingerprints yielded a mean $R^2$ of \MorganRTwoMean, without a significant difference from the 3D hybrid space ($R^2 = \RdkitTwoDThreeDFpRTwoMean$; Nadeau-Bengio-corrected paired $t$-test on the per-fold $R^2$, $p = \NBpRTwoMorganVsTwoDThreeDFp$). The two best-scoring spaces, the 2D descriptors with fingerprints and the 3D hybrid, were likewise indistinguishable from each other ($p = \NBpRTwoTwoDFpVsTwoDThreeDFp$). Since no space was reliably superior and the Morgan space needs neither descriptor calculation nor conformer generation, it was adopted for all downstream feature selection, data augmentation, and virtual screening protocols.


% ==========================================
% 2. SECTION 2.5
% ==========================================
\subsection{Nested Cross-Validation Protocol and Variance Correction}
\label{sec:materials_ncv}

Model building and model assessment were separated by a $\OuterRepeats \times \OuterKFolds$ repeated nested
cross-validation with random seed 42: a \OuterKFolds-fold outer loop repeated \OuterRepeats\ times, giving
\TotalOuterFolds\ outer evaluations, with a \InnerKFolds-fold inner loop inside each outer training partition
\citep{krstajic2014, cawley2010}. Methodological frameworks do not require the two loops to use the same number of
folds \citep{krstajic2014}, and a systematic comparison of validation schemes finds no uniformly superior variant, its
ranking depending on the algorithm, the data and the circumstances of the study \citep{heberger2023}. The same outer
partitions were used for every representation and every model, so that all comparisons reported here are paired.
Figure~\ref{fig:nested_cv} shows one outer fold of this protocol.

% Figure 2: one outer fold of the protocol of Section 2.4. The point of the figure is the hatched marker: every
% data-dependent transformation is estimated inside the partition that encloses it, never on the data it is scored on.
\begin{figure}[!h]
	\centering
	\resizebox{0.93\textwidth}{!}{%
		\begin{tikzpicture}[
			font=\sffamily\small,
			>=Stealth,
			node distance=6mm and 9mm,
			cvbox/.style={draw=gray!60!black, fill=gray!4, thick, rounded corners=3pt,
				align=center, text width=4.5cm, inner sep=5pt},
			outerstep/.style={cvbox, draw=blue!70!black, fill=blue!5},
			innerstep/.style={cvbox, draw=teal!70!black, fill=teal!5},
			fitstep/.style={cvbox, draw=orange!85!black, fill=orange!8, text width=4.9cm},
			testbox/.style={cvbox, draw=red!65!black, fill=red!4, text width=3.9cm},
			flow/.style={->, thick, draw=gray!65},
			note/.style={font=\sffamily\scriptsize\itshape, text=gray!45!black, align=left, text width=4.3cm}
			]

			\node[cvbox, text width=5.2cm] (data)
			{\textbf{Data set}\\$N = \NCompounds$ compounds};

			\node[outerstep, below=of data, text width=5.2cm] (split)
			{\textbf{Outer split}\\$\OuterKFolds$ folds $\times\ \OuterRepeats$ repeats $=\ \TotalOuterFolds$ evaluations\\\footnotesize shared by every representation and model};

			\node[outerstep, below left=8mm and 6mm of split] (pool)
			{\textbf{Outer training pool}};

			\node[testbox, below right=8mm and 6mm of split] (test)
			{\textbf{Outer test fold}\\\footnotesize held out from everything below};

			\draw[flow] (data) -- (split);
			\draw[flow] (split) -- (pool);
			\draw[flow] (split) -- (test);

			% ---- inner loop ----
			\node[innerstep, below=9mm of pool] (inner)
			{\textbf{Inner loop: $\InnerKFolds$ folds}\\\footnotesize $\NArchitectureCombos$ candidate architectures\\\footnotesize 1--$\MaxComboSize$ of RF, ET, LGBM, XGB, SVR, kNN};

			\node[fitstep, below=of inner] (prep)
			{\textbf{Preprocessing, fitted here}\\\footnotesize median imputation $\rightarrow$ scaling $\rightarrow$ feature selection\\\footnotesize estimated on the inner training fold only};

			\node[innerstep, below=of prep] (fit)
			{\textbf{Fit candidate}\\\footnotesize single model, or stacking with a RidgeCV meta-learner on out-of-fold predictions};

			\node[innerstep, below=of fit] (score)
			{\textbf{Score inner validation fold}};

			\node[innerstep, below=of score] (select)
			{\textbf{Select}\\\footnotesize highest mean inner $R^2$; hyperparameters fixed, so only the architecture is chosen};

			\draw[flow] (pool) -- (inner);
			\draw[flow] (inner) -- (prep);
			\draw[flow] (prep) -- (fit);
			\draw[flow] (fit) -- (score);
			\draw[flow] (score) -- (select);

			% ---- refit and outer scoring ----
			\node[fitstep, below=11mm of select] (refit)
			{\textbf{Refit on the whole outer training pool}\\\footnotesize the preprocessing is estimated again, on this pool};

			\node[outerstep, below=of refit, text width=4.9cm] (predict)
			{\textbf{Predict the outer test fold once}\\\footnotesize predictions clipped to the training range $\pm\ 3$ $pIC_{50}$};

			\node[outerstep, below=of predict, text width=5.6cm] (agg)
			{\textbf{Aggregate the \TotalOuterFolds\ outer folds}\\\footnotesize $R^2$ and MAE; Nadeau--Bengio corrected paired $t$-test, Holm adjustment, TOST at $\pm 0.05$};

			\draw[flow] (select) -- (refit);
			\draw[flow] (refit) -- (predict);
			\draw[flow] (predict) -- (agg);
			\draw[flow] (test.south) |- ($(predict.north)+(0,4mm)$) -| (predict.east);

			\begin{scope}[on background layer]
				\node[draw=teal!70!black, fill=teal!3, rounded corners=5pt, inner sep=7pt,
				fit=(inner)(prep)(fit)(score)(select),
				label={[text=teal!60!black, font=\sffamily\bfseries\small]above:Inner loop: architecture selection}] {};
				\node[draw=blue!70!black, fill=blue!2, rounded corners=5pt, inner sep=11pt,
				fit=(split)(pool)(test)(inner)(prep)(fit)(score)(select)(refit)(predict)(agg),
				label={[text=blue!60!black, font=\sffamily\bfseries\small]above:Outer loop: performance estimation}] {};
			\end{scope}

		\end{tikzpicture}%
	}
	\caption{\textbf{One outer fold of the evaluation protocol.} The outer loop estimates performance and the inner
		loop selects the architecture, so that no model is scored on the data that chose it. The two orange steps are
		the reason the estimate is leakage-free: imputation, scaling and supervised feature selection are not computed
		once on the whole data set but estimated again inside each partition they are used in, first on the inner
		training fold during selection and then on the outer training pool for the final fit, and applied to the
		held-out data without re-estimation. The outer test fold enters only at the prediction step. Because the
		\TotalOuterFolds\ outer partitions are shared by every representation and every model, all comparisons in this
		study are paired on the same molecules.}
	\label{fig:nested_cv}
\end{figure}


Random splitting was adopted for the main benchmark. The objective is interpolation within the applicability domain of
known TgDHFR inhibitors rather than extrapolation to new chemotypes, and at this sample size ($N = \NCompounds$)
scaffold-based partitioning produces imbalanced folds from which whole bioactivity ranges can disappear
\citep{kyrali2022}. Because random splits are optimistic for series of close analogs \citep{yang2019, aga2025}, the
Morgan-fingerprint ensemble and the D-MPNN were also run under scaffold-grouped splits. The Bemis-Murcko scaffold of
each compound, without stereochemistry, defined a group; the \NCompounds\ compounds fall into \ScaffoldNGroups\ groups,
the largest with \ScaffoldLargestGroup\ members and \ScaffoldSingletonGroups\ containing a single compound. In each of
the \OuterRepeats\ repeats, groups were visited in random order and each was assigned to the currently smallest of the
\OuterKFolds\ folds, so that no scaffold appears on both sides of an outer split. Everything else was identical to the
random-split protocol, and both model families again shared the same folds.

The inner loop selected the architecture and, at the same time, confined every data-dependent transformation to the
partition it was fitted on. Within each outer training pool it evaluated all \NArchitectureCombos\ single, pairwise and
three-way combinations of six regressor families (Random Forest, Extra Trees, LightGBM, XGBoost, SVR with an RBF kernel
and $C = 10$, and $k$-nearest neighbours with $k = 5$ and cosine distance; the tree ensembles used 200 trees) and kept
the combination with the highest mean inner-fold $R^2$. All hyperparameters were fixed in advance, so the inner loop
selects among architectures only; multi-model combinations were combined by a \texttt{StackingRegressor} with a
\texttt{RidgeCV} meta-estimator trained on 5-fold out-of-fold base predictions. Fitting a global transformation before
partitioning would leak data-set-wide statistics into the training phase \citep{cawley2010, varma2006}, so imputation,
standardization and supervised feature selection were implemented as steps of the model pipeline and their parameters
estimated only on the training partition at hand, an inner training fold during selection and the outer training pool
for the final fit, then applied to the held-out data without re-estimation \citep{brva2026}. For the continuous
descriptor spaces that pipeline chained median imputation, \texttt{StandardScaler} and \texttt{SelectFromModel} keeping
every descriptor with a strictly positive Random Forest importance; the binary Morgan space was passed to the
regressors unchanged. The selected combination was refit on the whole outer training pool and scored once on the
held-out outer fold, with predictions clipped to the training range extended by 3 $pIC_{50}$ units as a safeguard.

Repeated cross-validation violates the independence assumption of standard parametric tests, because its training sets
overlap across iterations. Every comparison between models was therefore made on the shared outer folds with the
Nadeau--Bengio corrected paired $t$-test \citep{nadeau2003}, which inflates the variance of the mean difference over the
$J = \TotalOuterFolds$ outer folds from $\hat{\sigma}^2/J$ to $(1/J + n_{test}/n_{train})\,\hat{\sigma}^2$, with
$n_{test}/n_{train} = 1/4$ for the \OuterKFolds-fold outer loop; $p$ values within each family of comparisons were then
adjusted by Holm's step-down procedure \citep{holm1979}. Because a non-significant difference is not evidence that two
models perform alike, each comparison was additionally submitted to two one-sided tests \citep{schuirmann1987} against
an equivalence margin of $\pm 0.05$ in $R^2$ and in MAE, equivalence being declared only when the 90\% confidence
interval of the difference lies entirely inside that margin. The margin is our own choice, as the literature offers no
standard value for it, and a wider one would make equivalence easier to declare. Confidence intervals for the hold-out
ablation of Section~\ref{sec:res_classical} come instead from 2000 percentile bootstrap resamples of the held-out
molecules \citep{efron1993}; the resampled indices depend only on the seed, so the same molecules are drawn for every
model and the differences are paired.

\subsection{Deep learning baseline: Directed message passing neural networks}
\label{sec:materials_gnn}

Benchmarking classical machine learning ensembles against advanced representation learning required implementing a Directed Message Passing Neural Network (D-MPNN) via the Chemprop framework (version 2.2.4) \citep{graff2026, yang2019}. To guarantee methodological parity, the D-MPNN was evaluated on the same \TotalOuterFolds\ outer partitions (\OuterRepeats$\times$\OuterKFolds\ repeated cross-validation, random seed 42) used for the classical baselines, with each molecule predicted by a model that never saw it. Two variants were run: a default configuration, for which no inner loop was needed, and a tuned variant with a hyperparameter search nested inside each outer fold. This alignment keeps the performance metrics directly comparable and removes the variance of a single split.

The architecture processes molecular SMILES strings by mapping them into graph representations, designating atoms as vertices and bonds as edges \citep{heid2024}. Initial feature vectors capturing atomic identity and topological connectivity were constructed using \texttt{SimpleMoleculeMolGraphFeaturizer()}. The network applied \texttt{BondMessagePassing()} (hidden size 300, depth 3, no dropout) to extract embeddings directly from directed edges rather than vertices. A \texttt{MeanAggregation()} function pooled these atomic embeddings into a unified molecular vector.

Within each training partition, the continuous $pIC_{50}$ targets were standardized with the training mean and standard deviation, and predictions were transformed back to the original scale. The pooled vector was batch-normalized and fed into a \texttt{RegressionFFN()} (one hidden layer of 300 units) to yield the final predictive output. The mean squared error on the standardized targets was the training loss. Model weights were optimized via PyTorch Lightning for a fixed number of 100 epochs on CPU, using Chemprop's default learning-rate schedule (2 warm-up epochs, initial rate $10^{-4}$, peak $10^{-3}$, final $10^{-4}$) and a per-fold random seed (42 plus the fold index). No early stopping or checkpoint selection was applied, so the model after the last epoch was used and the held-out fold never influenced training.

Unlike the classical ensembles, whose architecture was selected by an inner search, the default D-MPNN used Chemprop's default hyperparameters (hidden size 300, depth 3, no dropout). A search matching the classical one would require \GnnSearchFits\ trainings per outer fold (\NArchitectureCombos\ architecture combinations across \InnerKFolds\ inner splits), which at the measured cost of one training is not feasible for the full protocol (Section~\ref{sec:cost_results}); the tuned variant therefore uses a reduced search.

The tuned variant searched a grid of \GnnTunedNConfigs\ configurations (hidden size 300 or 600, message-passing depth 3 or 5, dropout 0 or 0.2; the default configuration is one of them). Inside each outer training pool, 20\% of the compounds were held out at random, every configuration was trained on the remaining 80\% and the one with the lowest root-mean-square error on the held-out part was retained. That configuration was then retrained on the whole outer training pool, with the same 100-epoch schedule and per-fold seed, and scored once on the outer test fold. The held-out outer fold therefore played no role in the choice. The search is deliberately smaller than the classical one (a single inner split instead of \InnerKFolds\ folds, and \GnnTunedNConfigs\ configurations instead of \NArchitectureCombos\ combinations), so it tests whether a moderate amount of tuning helps rather than the best result that exhaustive tuning could reach. The default variant separates the inductive bias of the graph architecture from the effect of tuning on a small sample \citep{cawley2010, hadas2023}, and the tuned variant checks whether that choice handicapped the network.

The default and the tuned variant both test one architecture on a modest budget. To separate what the graph
architecture contributes from what the training budget and external data contribute, four further neural models were
evaluated on the same \TotalOuterFolds\ outer folds, each fitted on the outer training partition alone and scored
exactly like the models above. The first is a feed-forward network on the same \ExpBNFeaturesTotal-bit ECFP4
fingerprints used by the classical ensembles, which keeps their input and learns its own transformation of it. The
second concatenates the vector learned by the D-MPNN with the RDKit 2D descriptors before the feed-forward head, so
that the network receives the descriptor block as well. The third averages the predictions of five D-MPNN that differ
only in their random seed, which buys accuracy with computation alone. The fourth initializes the message-passing
weights from CheMeleon \citep{burns2025}, a model pretrained on classical molecular descriptors computed for a large
unlabeled corpus, and is the only model in this study that sees information from outside the benchmark. None of the
four was tuned.

Finally, the architecture of the original study was re-evaluated without its leaks on the same \TotalOuterFolds\ outer
folds, in two variants: the deep ensemble on all descriptors (DNN baseline) and the complete original pipeline of
importance-based selection and Gaussian augmentation (DNN full), both with the selection and the augmentation confined
to each training partition and with early stopping on an internal split. This separates what the architecture
contributes from what the protocol contributes, and places the original model on the same folds as every other model
in this study.


% ==========================================
% 2. SECTION 2.7
% ==========================================
\subsection{Application: Screening within the Applicability Domain}
\label{sec:materials_screening}

The feature-selected Morgan-fingerprint stacking ensemble of the ablation above (Random Forest, XGBoost and SVR
combined by a ridge meta-learner, trained without augmentation on the \ExpNDev-compound development partition of the
15\% hold-out split; hold-out $R^2 = \ExpBRTwoMean$) was applied to a library of FDA-approved compounds
($N = \ScreenNLibrary$). Structural pre-filtering kept only compounds containing H, C, N, O, F, P, S, Cl, Br and I, a
set that covers all but \ScreenTrainNOutsideElements\ of the training compounds, and a size ceiling (molecular weight
$\le \ScreenMaxMW$ Da; the largest training compound weighs \ScreenTrainMaxMW\ Da) removed molecules whose potency the
model could not meaningfully extrapolate.

Prediction on a library of this kind is only meaningful where the model has support. Following the finding of Sheridan
et al.\ that prediction error grows as a molecule becomes less similar to the training set \citep{sheridan2004}, every
FDA compound was given its mean Tanimoto similarity (2048-bit Morgan fingerprints) to its \ScreenADNeighbors\ most
similar training compounds, and was flagged as inside the applicability domain when this value was at least
\ScreenADMinSimilarity, the threshold obtained for the training compounds themselves as the mean plus \ScreenADZ\
standard deviations of their distances \citep{sahigara2012}. The multiplier is an arbitrary choice: the literature
offers no universal rule, and a comparison of applicability-domain methods concludes that the appropriate one has to be
determined for the data set at hand \citep{heberger2023}. For that reason no compound was excluded on this flag, which
is reported with every candidate.

Candidates were ranked by predicted potency alone. Four ligand efficiency metrics were computed from the prediction and
reported beside it, to show how much of a predicted potency is accounted for by molecular size \citep{hopkins2004,
hopkins2014}; none of them was used as a filter, and their definitions and the assumptions they carry are given in the
Supplementary Material. The ranked candidates were then submitted to an orthogonal structure-based check with AutoDock
Vina against the TgDHFR crystal structure 6AOG \citep{trott2009, consortium2019}, whose protocol, validated by
re-docking the co-crystallized ligand to a heavy-atom RMSD of $\RedockRMSD\text{ \AA}$, is described in the
Supplementary Material together with the consensus thresholds.

This stage generates computational hypotheses. None of the candidates was tested experimentally, and selectivity
against human DHFR, which the antifolates among them are known to inhibit, was not assessed.

\subsection{Computational Environment and Reproducibility}

To support algorithmic transparency and computational reproducibility \citep{rojas2026}, all data preparation, model training, and evaluation steps were programmatically implemented in Python (version 3.11.16) and chained in a single entry point that runs every step of the study in order, from the construction of the training sets to the final consistency check of the manuscript. The computational environment relied on a standardized core software stack:

\begin{itemize}
	\item \textbf{Cheminformatics and Featurization:} RDKit (version 2025.9.3) \citep{landrum2025} executed all molecular parsing, structural standardization, and descriptor generation (2D, 3D, and ECFP4); Open Babel (version 3.1.1) \citep{babel2011} converted the receptor to PDBQT.
	\item \textbf{Classical Machine Learning:} The nested cross-validation framework \citep{krstajic2014, calle2025} and baseline algorithms leveraged \texttt{scikit-learn} (version 1.6.0) \citep{pedregosa2011sklearn}, \texttt{xgboost} (version 3.1.2) \citep{chengeustrin2016XGBoost}, and \texttt{lightgbm} (version 4.6.0) \citep{kemeng2017LightGBM}. Data manipulation, matrix operations, and statistical tests were managed via \texttt{pandas} (version 2.3.3) \citep{mckinney2010pandas}, \texttt{numpy} (version 1.26.4) \citep{harris2020numpy}, and \texttt{scipy} (version 1.16.3).
	\item \textbf{Deep Learning Infrastructure:} The D-MPNN graph models utilized \texttt{chemprop} (version 2.2.4) \citep{graff2026, heid2024} and \texttt{lightning} (version 2.6.5), with underlying tensor operations supported by PyTorch (\texttt{torch}, version 2.5.1) \citep{paszke2019pytorch}.
	\item \textbf{Docking:} AutoDock Vina (Python bindings, version 1.2.7) \citep{trott2009} performed the docking and ligands were prepared with Meeko (version 0.8.0).
	\item \textbf{Hardware Specifications:} All model training, architecture selection, and evaluation protocols ran on an Ubuntu 26.04 LTS workstation with two Intel(R) Xeon(R) E5-2695 v2 processors (24 physical cores, 48 threads, 2.4 GHz) and 247 GB of RAM, using CPU only; no GPU was involved.
\end{itemize}

All random number generators were seeded (seed 42). The models of the augmentation study were fitted single-threaded, because a multi-threaded random forest sums its trees in a thread-dependent order and changes the last bits of its predictions; parallelism was applied across independent tasks (features in the permutation importance, runs in the ensemble) so that the results do not depend on the number of workers. Together with the numerical tolerance of the selection rule (Section~\ref{sec:materials_curation}), this made the feature selection reproducible: permutation importances computed with 46 and with 8 workers were bitwise identical. The graph-model benchmark, in contrast, is not bit-reproducible: training the network on several CPU threads sums floating-point terms in a thread-dependent order, which a fixed seed does not control. Three executions of the same code on the same outer folds differed by 0.004 in mean fold $R^2$, well inside the confidence interval of the estimate itself ([\ChemPropDMPNNFoldCILow, \ChemPropDMPNNFoldCIHigh]), so the variation affects the third decimal of the neural results and none of the comparisons reported here. Every classical result, by contrast, reproduces exactly, because the outer folds come from the fixed seed and both the preprocessing and the architecture search are confined to each training fold. The wall-clock time of every step was recorded on this machine, without normalizing by thread count, and is reported in Section~\ref{sec:cost_results}.
